%% appendix_bayes_pedagogy.tex — Bayesian Inference Background
%% Explanatory material moved from Methods chapter.

\chapter{Bayesian Inference: Background and Derivation}
\label{app:bayes_pedagogy}

This appendix provides the conceptual background for the Beta-conjugate
Bayesian framework used in Chapter~\ref{ch:methods}. Readers familiar
with conjugate Bayesian updating may skip this appendix.

\section{The Latent Variable}
\label{sec:latent_variable}

In this framework, habitability $H$ is treated as a \textit{latent variable}:
a quantity that cannot be directly observed by any Cassini instrument but
whose value is inferred from the combination of observational proxies
$\featurevec = (f_1, \ldots, f_8)$.  A latent variable (from the Latin
\textit{latere}: to be hidden) influences observable outcomes but is not
itself directly measurable.  Here, $H$ is not a binary yes-or-no property
but a continuous probability in $[0, 1]$: the value $H = p$ at a given pixel
means that pixel has probability $p$ of satisfying the relevant habitability
conditions given all available evidence.  Each feature $f_i$ provides partial,
indirect evidence about the value of $H$.

\section{Why a Beta Distribution?}
\label{sec:why_beta}

The Beta distribution $\mathrm{Beta}(\alpha, \beta)$ is defined on $[0, 1]$
and characterised by two positive shape parameters $\alpha$ and $\beta$.
Its probability density function (PDF) is
\begin{equation}
  p(h) = \frac{\Gamma(\alpha + \beta)}{\Gamma(\alpha)\,\Gamma(\beta)}\;
  h^{\,\alpha - 1}\,(1 - h)^{\,\beta - 1}, \qquad h \in [0,1]
  \label{eq:prior_pdf}
\end{equation}
where $\Gamma(\cdot)$ is the gamma function, which generalises the factorial to
non-integer arguments and is defined for all positive reals by
\begin{equation}
  \Gamma(x) = \int_0^{\infty} t^{\,x-1}\,e^{-t}\,\mathrm{d}t .
  \label{eq:gamma_def}
\end{equation}
For positive integers it reduces to the factorial, $\Gamma(n) = (n-1)!$, and it
obeys the recurrence $\Gamma(x+1) = x\,\Gamma(x)$ (the continuous generalisation
of $n! = n\,(n-1)!$); for non-integer arguments such as $\alpha_0 = 1.655$ it is
evaluated numerically (here via \texttt{scipy.special.gamma}).  The ratio of
gamma functions in Eq.~\ref{eq:prior_pdf} is the reciprocal of the Beta function
$\mathrm{B}(\alpha,\beta) = \Gamma(\alpha)\,\Gamma(\beta)/\Gamma(\alpha+\beta)$,
the normalising constant that makes $p(h)$ integrate to one.  For the prior
adopted in Chapter~\ref{ch:methods} ($\alpha_0 = 1.655$, $\beta_0 = 3.345$ from
$\kappa = 5$, $\mu_0 = 0.331$), this evaluates to
\begin{equation}
  p_0(h) = \frac{\Gamma(5)}{\Gamma(1.655)\,\Gamma(3.345)}\;
  h^{\,0.655}\,(1 - h)^{\,2.345}
  \approx 9.47\;h^{\,0.655}\,(1 - h)^{\,2.345} .
  \label{eq:prior_pdf_numeric}
\end{equation}
The distribution is unimodal when both $\alpha > 1$ and $\beta > 1$, with mean
$\alpha/(\alpha + \beta)$ (here $\mu_0 = 0.331$) and mode
$(\alpha-1)/(\alpha+\beta-2) \approx 0.218$.  Its standard deviation follows from
the Beta variance $\mathrm{Var} = \alpha\beta/[(\alpha+\beta)^2(\alpha+\beta+1)]$
\citep[p.~581, Table~A.1]{Gelman2013}:
\begin{equation}
  \sigma_0 = \sqrt{\frac{\alpha_0\,\beta_0}{(\alpha_0+\beta_0)^2\,(\alpha_0+\beta_0+1)}}
  = \sqrt{\frac{1.655 \times 3.345}{25 \times 6}} \approx 0.19 .
  \label{eq:prior_sd}
\end{equation}
Because it is bounded on $[0,1]$ and can represent any unimodal belief over a
probability, the Beta distribution is the natural prior for the latent
habitability probability $H$.

\section{Bernoulli Trials and the Binomial Likelihood}
\label{sec:bernoulli}

A \textit{Bernoulli trial} is an experiment with exactly two possible outcomes (conventionally called ``success'' and ``failure'') each occurring with fixed probabilities $p$ and $1-p$ respectively \citep[ch.~2]{DeGroot2012}.  A coin flip is the prototypical example: heads (success, probability $p$) or tails (failure, probability $1-p$).  The \textit{Binomial distribution} counts the total number of successes $k$ in $n$ independent Bernoulli trials; it is the natural likelihood function for data arising from repeated binary observations.

In the habitability model, no actual coin flips occur.  Instead, the prior and data contributions are expressed as \textit{virtual} (or \textit{pseudo}) Bernoulli trials: a mathematical device for encoding strength of belief in the same currency as observational data.  The prior contributes $\kappa = 5$ virtual trials, of which $\alpha_0 = 1.655$ are ``successes'' (prior evidence for habitability) and $\beta_0 = 3.345$ are ``failures'' (prior evidence against).  Each observed feature $f_i$ then contributes $\lambda w_i$ additional virtual trials, of which a fraction $f_i$ are successes and $(1-f_i)$ are failures.  This pseudo-count interpretation is the standard Bayesian device for placing prior beliefs on the same mathematical footing as observed data \citep[ch.~2]{Gelman2013}, and is the mechanism by which conjugacy yields a closed-form posterior (Section~\ref{sec:conjugacy_explained}).

\section{Conjugacy}
\label{sec:conjugacy_explained}

A prior distribution is \textit{conjugate} to a likelihood function if the
posterior (prior updated by data) has the same mathematical form as the
prior, differing only in its parameter values \citep[ch.~2]{Gelman2013}.  The Beta distribution is
conjugate to the Binomial likelihood: if the prior on a success probability
$p$ is $\mathrm{Beta}(\alpha_0, \beta_0)$ and we observe $k$ successes
in $n$ Bernoulli trials, the posterior is
$\mathrm{Beta}(\alpha_0 + k,\; \beta_0 + (n - k))$.
In the habitability model, this update takes the specific form given in
Eqs.~\ref{eq:alpha_post}--\ref{eq:beta_post}, where the pseudo-counts
$k = \lambda\sum_i w_i f_i$ and $n - k = \lambda\sum_i w_i(1-f_i)$ are
derived from the weighted feature values.  The resulting posterior PDF has
the same gamma-function structure as the prior PDF (Eq.~\ref{eq:prior_pdf}),
with only the shape parameters changed:
\begin{equation}
  p\bigl(h \mid \featurevec\bigr)
  = \frac{\Gamma(\alpha_\text{post} + \beta_\text{post})}{\Gamma(\alpha_\text{post})\,\Gamma(\beta_\text{post})}\;
  h^{\,\alpha_\text{post} - 1}\,(1 - h)^{\,\beta_\text{post} - 1} ,
  \label{eq:posterior_pdf}
\end{equation}
so no new functional form is introduced.  Here $\alpha_\text{post} +
\beta_\text{post} = \kappa + \lambda = 11$ at every pixel, so the leading
factor is the constant $\Gamma(11) = 10! = 3{,}628{,}800$.
This analytical tractability makes the update feasible at 6.49 million
pixels simultaneously with no numerical integration required.

\section{The Beta-Binomial Model Applied to Habitability}
\label{sec:beta_binom_detail}

Each feature $f_i \in [0, 1]$ is treated as the fractional outcome of
$\lambda w_i$ virtual Bernoulli trials: fraction $f_i$ of those trials are
``successes'' (evidence for habitability) and $(1 - f_i)$ are ``failures''
(evidence against).  High feature values contribute more positive
pseudo-observations and shift the posterior towards higher habitability;
low values do the opposite.  The prior contributes $\kappa$ virtual
observations and the data from all eight features contribute a further
$\lambda \sum_i w_i = \lambda$ observations.

Figure~\ref{fig:beta_update} visualises the updating process for three
representative sites (Towada Lacus, Belet dune field, Selk crater),
showing how each successive feature narrows and displaces the Beta
distribution from the prior to the final posterior.  The figure also
shows the 95\%~CI for eight key present-epoch sites.

Figure~\ref{fig:bayes_update} illustrates this update graphically for three
representative sites.  The left panel shows the prior
$\mathrm{Beta}(\alpha_0, \beta_0)$, which is weakly informative
($\sigma \approx 0.19$) and centred at $\mu_0 = 0.331$.  The right panel
shows the posteriors after incorporating the respective Cassini feature
values: Towada Lacus shifts strongly rightward
(highest VIMS organic abundance combined with high liquid and methane-cycle scores), while the Belet dune
sea shifts only modestly (moderate organic but no liquid) and Selk crater shifts
leftward (below-average methane cycle and surface--atmosphere scores despite
exceptional geodiversity).